\documentclass[11pt]{letter}

\usepackage[utf8]{inputenc}
\usepackage[english]{babel}
\usepackage[margin=1in]{geometry}
\usepackage{parskip}
\usepackage{hyperref}

\begin{document}

Dear \textit{eLife} Editors,

We submit our manuscript, \textit{Living Alignment}, for consideration as a Review Article at \textit{eLife}. 

The paper first highlights an observation that has been made independently in several disciplines: living systems are healthy and flourishing when they avoid collapse to patterns of organization that admit excessively simple descriptions. We make this principle concrete with examples from a range of scales -- evolutionary biology, cell biology, neuroscience, psychology and cognitive science, sociology, ecology.

We then apply this insight from the life sciences to systematize 

Although researchers disagree on details of how aligned AI systems should behave, there is broad agreement that the goal of alignment is futures full of health and flourishing. We argue this means that 


understanding what health looks like in living systems is central to the alignment problem. 

We then argue that if the AI alignment problem is 

We then argue that the AI alignment problem cannot be separated from 

The problem of AI alignment in a very broad sense -- what kinds of futures are full of health and flourishing -- is actually a question about life.

Studying life itself, in the broadest sense, grounds how we can think about AI and alignment.

Systems with rich internal diversity and structured interactions 

create potential for fundamentally new patterns of organization to emerge over time.

We link this principle between .

We then 

Our premise is that making good AI can't be separated from understanding life. 

AI comes into existence in a living world. 

A deep understanding of life is necessary to build aligned AI systems.


For example, at some point in the history of life, DNA came into existence. Later, multicellular life came into existence.

This hypothesis leans on systems science and physics perspectives about the dynamics of living systems and extends it to capture phenomena at many levels of scale.


can be taken as a broad perspective on life. But because life is 



The review asks what living systems---from cells to ecosystems to societies---can teach us about AI alignment, and proposes that healthy living systems share a common principle: they avoid \textit{collapse} to any single form by maintaining intricate networks of semi-permeable boundaries that hold diverse, partial perspectives in generative interplay.

\textbf{How the review moves the field forward.} Most existing alignment writing operates inside one of two frames: technical (specifying objectives, evaluating models, interpreting circuits) or normative (whose values, which principles). Our review offers a third frame---a systems-level, biologically grounded account of what alignment is \textit{for}---and shows that this frame is not merely metaphorical. We demonstrate that several canonical alignment failures (perverse instantiation, instruction-following failures, conceptual monoculture, belief reinforcement, inequality) can be re-described as instances of a single underlying failure mode: collapse of contextual sensitivity to a myopic form. We then argue that this re-description is generative: it predicts that alignment cannot be fully formalised, it clarifies why semi-permeable boundaries (rather than fixed objectives or fixed values) are the appropriate mechanism for keeping AI aligned as capabilities scale, and it gives concrete shape to what an aligned long-term future might look like even under radical changes to human society and to AI itself. For \textit{eLife}'s readership, the review also makes the reverse contribution: by putting alignment in dialogue with cellular biology, ecology, neuroscience, and group cognition, it surfaces a shared structural problem---the avoidance of collapse---that we believe will be of interest to life scientists studying any of these systems in their own right.

\textbf{Why we are suitable authors.} We combine expertise across several disciplines to provide a broad perspective on life and alignment. ZKN, MMN and MGM are computational neuroscientists and psychiatrists who study how biological brains maintain flexible representations and how those representations collapse in illness. JZL and NT are senior researchers at Google DeepMind working on cooperative AI, multi-agent systems and AI safety. SS works at the interface of clinical science and human flourishing at OHSU, and EM has spent over a decade applying behavioural and cognitive science to consumer-facing technology. This combination of biological, clinical, computational and applied AI perspectives is, we believe, exactly what is required to write a review that takes life seriously as a source of insight for alignment. Our group has also engaged with this specific question before: a subset of the authors developed the related ``dynamic diversity'' framing for proxy failure (Kurth-Nelson et al., 2024), and the present manuscript substantially generalises and extends that earlier line of thinking.


\textbf{How it differs from recent reviews in the field.} 

Different from broad perspectives on living systems (what is life, friston, allostasis, autopoeisis, etc). 

Different from reviews on alignment (gabriel, sorensen, leibo, hannah rose kirk).

And the synthesis of the two.

Recent reviews of AI alignment have made important contributions but address different questions. Anwar et al.'s \textit{Foundational Challenges in Assuring Alignment and Safety of Large Language Models} (2024) is a comprehensive technical catalogue of open problems specific to current LLM systems; our review is largely orthogonal to it, asking instead what general principles should govern alignment across systems and across time. Gabriel's \textit{Artificial Intelligence, Values, and Alignment} (Minds and Machines, 2020) and Gabriel \& Keeling's \textit{A Matter of Principle?} (Philosophical Studies, 2025) approach alignment as a normative problem in political philosophy---how to choose principles fairly---whereas we approach it as a problem in the dynamics of living systems and argue that no fixed principle can suffice. Sorensen et al.'s \textit{A Roadmap to Pluralistic Alignment} (2024) tackles the related problem of accommodating diverse human values, but treats plurality primarily as a property of the target rather than, as we argue, a property of the entire human--AI--biosphere system that must be sustained dynamically. Bereska \& Gavves's \textit{Mechanistic Interpretability for AI Safety---A Review} (2024) and Kirk et al.'s review of personalised alignment (\textit{Nature Machine Intelligence}, 2024) review specific technical sub-areas. To our knowledge, no recent review situates AI alignment within a unifying framework drawn from how cells, brains, ecosystems and societies achieve and lose health---nor draws the connection through the concept of semi-permeable boundaries to the multi-scale machinery that prevents collapse in each of these systems.

We hope you find the manuscript a good fit for \textit{eLife}, and we look forward to your response. We have no conflicts of interest to disclose beyond those stated in the manuscript.

\closing{With kind regards,}

\end{document}
